\Article{13}
As our last example, we will consider the general representation of the surface of an ellipsoid of revolution on a spherical surface. For the ellipsoid we retain the notation of the previous article. With the radius of the sphere \(=A\), we set
\[
\begin{aligned}
X&=A\cos T\sin U\\
Y&=A\sin T\sin U\\
Z&=A\cos U
\end{aligned}
\]
\SourceLeaf{213}%
\Needspace{10\baselineskip}
\noindent Applying the general solution of Article 5 here, we find that, with \(f\) denoting an arbitrary function, \(T\) and \(i\log\mathrm{cotang}\tfrac{1}{2}U\) must be set equal to the real and imaginary parts of
\[
f(t+i\log\{\mathrm{cotang}\tfrac{1}{2}w\mathord{.}(\frac{1-\varepsilon\cos w}{1+\varepsilon\cos w})^{\frac{1}{2}\varepsilon}\})
\]
\noindent respectively\footnote{We pass over here both the second solution of Article 5, which would differ from the one above only by exchanging \(-T\) for \(+T\) and would give a reversed representation, and the case of an elongated ellipsoid, whose treatment follows directly from that of the flattened ellipsoid in view of what was set out in the previous article.}.

The simplest solution is to set \(f\upsilon=\upsilon\), giving
\[
T=t,\qquad \mathrm{tang}\tfrac{1}{2}U=\mathrm{tang}\tfrac{1}{2}w\mathord{.}(\frac{1+\varepsilon\cos w}{1-\varepsilon\cos w})^{\frac{1}{2}\varepsilon}
\]
\noindent This provides a transformation of great use in higher geodesy, though here we can give only a few brief indications of its use. Consider points on the ellipsoid and on the sphere to correspond if they have the same longitude and if their latitudes, \(90^0-w\) and \(90^0-U\) respectively, are related by the equation given above. Then a system of comparatively small triangles on the spheroid formed by shortest lines corresponds to a system of triangles on the sphere. Triangles usable for actual measurement will always be comparatively small. The angles of the spherical triangles are \textit{exactly} equal to the corresponding angles on the spheroid. Their sides depart so little from arcs of great circles that in most cases, where the utmost precision is not required, they can be regarded as coinciding with those arcs. Where the greatest accuracy is required, the departure from the great circle can readily be calculated to the necessary precision by simple formulas. Thus, after first transferring one side of a triangle properly to the sphere, we can calculate the whole system from the angles just as if the system lay on the sphere itself, with the modification just indicated if necessary. For every point in the system we can determine \(T\) and \(U\) and go back from the latter values to the corresponding values of \(w\), most easily with an auxiliary table that is very easy to construct.
